% ==============================================================================
% Section 5: Sensitivity Analysis: Hindustan Unilever Break-Even Dynamics
% ==============================================================================

\clearpage
\section[Sensitivity Analysis]{Sensitivity Analysis: Hindustan Unilever Break-Even Dynamics}

\subsection{Algebraic Derivation of the Break-Even Perpetual Growth Rate}
Because the Gordon Growth Model displays asymptotic convexity as $g \to \Ke$, the valuation of Hindustan Unilever (\inr 2,512 at $g = 12.0\%$ vs Market Price of \inr 1,973) is subjected to rigorous sensitivity analysis. We analytically solve for the \textbf{break-even growth rate ($\bm{g^*}$)} where intrinsic value exactly equates to the current market quotation:

\begin{equation}
    P_0 = \frac{\text{Proxy DPS} \times (1 + g)}{\Ke - g} \implies 1,973 = \frac{44.86\,(1 + g)}{0.14 - g}
\end{equation}

Expanding and solving for $g$:
\begin{align}
    1,973 \times (0.14 - g) = 44.86\,(1 + g) &\implies 276.22 - 1,973\,g = 44.86 + 44.86\,g \\
    231.36 = 2,017.86\,g &\implies g^* = \frac{231.36}{2,017.86} \approx \bm{11.46\%}
\end{align}

\begin{formulabox}[HUL Growth-Rate Decision Boundaries]
\centering\footnotesize
$\bm{g > 11.46\%}$: \textbf{Undervalued} \quad $\bm{\cdot}$ \quad
$\bm{g = 11.46\%}$: \textbf{Fairly Valued} \quad $\bm{\cdot}$ \quad
$\bm{g < 11.46\%}$: \textbf{Overvalued}
\end{formulabox}

\subsection{Graphic Visualization of DDM Growth Sensitivity}
Figure~\ref{fig:hul_sensitivity} plots HUL's DDM intrinsic value across a plausible perpetual growth domain ($8.0\%$ to $12.8\%$) against the current secondary market price of \inr 1,973.

\begin{figure}[H]
\centering
\begin{tikzpicture}
\begin{axis}[
    width=13.0cm,
    height=5.2cm,
    xmin=8.0, xmax=13.0,
    ymin=500, ymax=3500,
    xlabel={\small Sustainable Perpetual Growth Rate $g$ (\%)},
    ylabel={\small DDM Value (\inr)},
    xtick={8, 9, 10, 11, 11.46, 12, 13},
    xticklabels={8\%, 9\%, 10\%, 11\%, \textbf{11.46\%}, 12\%, 13\%},
    ytick={500, 1000, 1500, 1973, 2500, 3000, 3500},
    yticklabels={\inr 500, \inr 1000, \inr 1500, \textbf{\inr 1973}, \inr 2500, \inr 3000, \inr 3500},
    grid=both,
    grid style={line width=0.3pt, draw=lightgrayborder},
    axis line style={line width=0.8pt, draw=pureblack},
    tick label style={font=\scriptsize},
    label style={font=\small},
    legend pos=north west,
    legend style={font=\scriptsize, draw=pureblack, fill=white}
]

    % DDM Curve: P0(g) = 44.86*(1 + x/100) / (0.14 - x/100)
    \addplot[domain=8.0:12.6, samples=80, line width=1.3pt, color=pureblack] 
        {44.86 * (1 + x/100) / (0.14 - x/100)};
    \addlegendentry{DDM Curve $P_{\DDM}(g)$}

    % Market Price Reference Line
    \addplot[domain=8.0:13.0, dashed, line width=1.0pt, color=midgray] {1973};
    \addlegendentry{Market Price ($P_0 = \inr 1,973$)}

    % Vertical Break-Even Line
    \draw[dotted, line width=1.0pt, color=darkgraytext] (axis cs:11.4655, 500) -- (axis cs:11.4655, 1973);

    % Break-Even Critical Point
    \addplot[only marks, mark=*, mark size=2.5pt, color=pureblack] coordinates {(11.4655, 1973)};
    \node[anchor=south east, font=\scriptsize\bfseries] at (axis cs:11.42, 2040) {Break-Even: $g^* = 11.46\%$};

    % Regional Annotations
    \node[font=\scriptsize\bfseries, align=center] at (axis cs:8.8, 1600) {Overvalued Region\\($P_{\DDM} < P_0$)};
    \node[font=\scriptsize\bfseries, align=center] at (axis cs:11.6, 3200) {Undervalued Region\\($P_{\DDM} > P_0$)};

\end{axis}
\end{tikzpicture}
\caption{Hindustan Unilever: DDM Intrinsic Value Sensitivity to Growth Rate $g$ ($\Ke = 14.0\%$).}
\label{fig:hul_sensitivity}
\end{figure}

\subsection{Empirical Implications of Growth Sensitivity}
The empirical takeaway is that HUL's apparent undervaluation is exceptionally fragile: a minor deceleration of just \textbf{54 basis points} (from the assumed $12.00\%$ to $11.46\%$) completely eliminates the margin of safety, rendering the corporation overvalued under the Gordon Growth framework.
